\section{Dataset and Sequence Preprocessing}
\label{sec:preprocessing}

We used the NASA C-MAPSS FD001 data.  Each observation contains an engine
identifier, cycle, three operating settings, and 21 sensor measurements.
Seven constant or near-constant columns were removed, leaving the 14 sensors
used by the sequence models.  The training RUL target was
\[
y_{i,t}=\max_{u \in i}(u)-t,
\]
where $i$ denotes an engine and $t$ its current cycle.  The 100 validation
cutoffs were kept fixed and the split was performed by engine: 80 engines
were used for fitting and 20 for validation.  Consequently, no window
contained observations from both partitions.

The scaler was fitted only on observations from the 80 training engines.
For sensor $j$, the transformed value was
$z_{i,t,j}=(x_{i,t,j}-\mu_j)/\sigma_j$.  At cycle $t$, a window of length $w$
was constructed as
\[
X_{i,t}=[z_{i,t-w+1},\ldots,z_{i,t}],
\]
using only observations at or before $t$.  If fewer than $w$ observations
were available, the first observed vector was repeated on the left.  Thus
short test engines remain usable without look-ahead or deletion.  A cap can
be applied to training targets as $\tilde y=\min(y,c)$, but the default
pipeline leaves targets uncapped and every evaluation uses the true RUL.

Models were evaluated in the exact order of the 100 fixed validation points.
We reported mean absolute error (MAE) and root mean squared error (RMSE) in
cycles, together with early-life ($y>125$) and late-life ($y\leq125$)
breakdowns.  Test windows were generated only at each engine's final
observed cycle; their labels were not used during model development.
